% GENERATED FILE. Edit canonical lessons, then run npm run generate:books.
% canonical-source-hash: fnv1a64:8659b73db066b180
% canonical-lessons: ML-C540-klinikku, ML-C540-unnuvati, ML-C540-mukhakkuru, ML-C540-maravikkuka, ML-C540-kalaykkuka, ML-R540-review-533-540

\chapter{At the Clinic}
\label{ch:ml-at-the-clinic}
\hlchaptermodality{540}

\begin{chapteropening}
\textbf{By the end of this chapter:} \emph{I can say a clinic, a walking stick, a pimple, to go numb and to feel sore.}

Complete the last lesson of chapter 540: I can say a clinic, a walking stick, a pimple, to go numb and to feel sore.
\end{chapteropening}

\section[\texorpdfstring{\ml{ക്ലിനിക്ക്} (klinikkŭ)}{ക്ലിനിക്ക് (klinikkŭ)}]{\ml{ക്ലിനിക്ക്} (klinikkŭ) — a clinic}

\label{lesson:ML-C540-klinikku}



\begin{quote}
Before the new one: say the Malayalam for a diet, then the Malayalam for nutritious food.
\end{quote}

\subsection*{You'll want to know: \ml{ക്ലിനിക്ക്}}
\textbf{\ml{ക്ലിനിക്ക്}} — \emph{klinikkŭ} — \textquotedblleft{}a clinic\textquotedblright{}.

\begin{quote}
\ml{ഞങ്ങളുടെ} \ml{കവലയിൽ} \ml{ഒരു} \ml{ചെറിയ} \ml{ക്ലിനിക്ക്} \ml{ഉണ്ട്}.
\end{quote}

\begin{quote}
\emph{ñaṅṅaḷuṭe kavalayil oru ceṟiya klinikkŭ uṇṭŭ.}
\end{quote}

\begin{quote}
There is a small clinic at our junction.
\end{quote}

Smaller than an \ml{ആശുപത്രി}, often one doctor's rooms.

\begin{cousinweb}
\textbf{\ml{ക്ലിനിക്ക്}} (\emph{klinikkŭ}) was borrowed from English \emph{clinic}.
\end{cousinweb}

\begin{grammarlens}[title={What you've built}]
One more word for describing what you lived and felt.
\end{grammarlens}

\subsection*{Your turn}
\begin{itemize}
\raggedright
  \item \emph{Say it:} \emph{klinikkŭ}
  \item \emph{Say it:} \emph{klinikkŭ}, once more
  \item \emph{Say it:} \emph{ñaṅṅaḷuṭe kavalayil oru ceṟiya klinikkŭ uṇṭŭ.} — There is a small clinic at our junction.
  \item \emph{Recall:} say the Malayalam for a diet, then the Malayalam for nutritious food, then say \emph{klinikkŭ} again
\end{itemize}

\begin{tcolorbox}[breakable,colback=teal!4,colframe=teal!35!black,title={Before you move on}]
What does \ml{ക്ലിനിക്ക്} mean? (\textquotedblleft{}A clinic\textquotedblright{}.) Say it once more.
\end{tcolorbox}
\section[\texorpdfstring{\ml{ഊന്നുവടി} (ūnnuvaṭi)}{ഊന്നുവടി (ūnnuvaṭi)}]{\ml{ഊന്നുവടി} (ūnnuvaṭi) — a walking stick; a crutch}

\label{lesson:ML-C540-unnuvati}



\begin{quote}
Before the new one: say the Malayalam for nutritious food, then the Malayalam for a clinic.
\end{quote}

\subsection*{You'll want to know: \ml{ഊന്നുവടി}}
\textbf{\ml{ഊന്നുവടി}} — \emph{ūnnuvaṭi} — \textquotedblleft{}a walking stick; a crutch\textquotedblright{}.

\begin{quote}
\ml{കാൽ} \ml{ഒടിഞ്ഞ} \ml{ശേഷം} \ml{അവൻ} \ml{ഊന്നുവടിയുമായി} \ml{നടന്നു}.
\end{quote}

\begin{quote}
\emph{kāl oṭiñña śēṣaṁ avan ūnnuvaṭiyumāyi naṭannu.}
\end{quote}

\begin{quote}
After breaking his leg, he walked with a crutch.
\end{quote}

\ml{ഊന്നുക} to lean on + \ml{വടി} a stick.

\begin{cousinweb}
\textbf{\ml{ഊന്നുവടി}} (\emph{ūnnuvaṭi}) is native Dravidian, a compound of \ml{ഊന്നുക}, \textquotedblleft{}to lean, to rest upon\textquotedblright{} (DEDR 763), and \ml{വടി}, \textquotedblleft{}a stick\textquotedblright{} (DEDR 5224).
\end{cousinweb}

\begin{grammarlens}[title={What you've built}]
One more word for describing what you lived and felt.
\end{grammarlens}

\subsection*{Your turn}
\begin{itemize}
\raggedright
  \item \emph{Say it:} \emph{ūnnuvaṭi}
  \item \emph{Say it:} \emph{ūnnuvaṭi}, once more
  \item \emph{Say it:} \emph{kāl oṭiñña śēṣaṁ avan ūnnuvaṭiyumāyi naṭannu.} — After breaking his leg, he walked with a crutch.
  \item \emph{Recall:} say the Malayalam for nutritious food, then the Malayalam for a clinic, then say \emph{ūnnuvaṭi} again
\end{itemize}

\begin{tcolorbox}[breakable,colback=teal!4,colframe=teal!35!black,title={Before you move on}]
What does \ml{ഊന്നുവടി} mean? (\textquotedblleft{}A walking stick; a crutch\textquotedblright{}.) Say it once more.
\end{tcolorbox}
\section[\texorpdfstring{\ml{മുഖക്കുരു} (mukhakkuru)}{മുഖക്കുരു (mukhakkuru)}]{\ml{മുഖക്കുരു} (mukhakkuru) — a pimple; acne}

\label{lesson:ML-C540-mukhakkuru}



\begin{quote}
Before the new one: say the Malayalam for a clinic, then the Malayalam for a walking stick.
\end{quote}

\subsection*{You'll want to know: \ml{മുഖക്കുരു}}
\textbf{\ml{മുഖക്കുരു}} — \emph{mukhakkuru} — \textquotedblleft{}a pimple; acne\textquotedblright{}.

\begin{quote}
\ml{അനിയത്തിക്ക്} \ml{മുഖത്ത്} \ml{ഒരുപാട്} \ml{മുഖക്കുരു} \ml{ഉണ്ട്}.
\end{quote}

\begin{quote}
\emph{aniyattikkŭ mukhattŭ orupāṭŭ mukhakkuru uṇṭŭ.}
\end{quote}

\begin{quote}
My little sister has a lot of pimples on her face.
\end{quote}

\ml{മുഖം} face + \ml{കുരു} a pimple, a boil.

\begin{cousinweb}
\textbf{\ml{മുഖക്കുരു}} (\emph{mukhakkuru}) is a Malayalam compound built from \ml{മുഖം}, from Sanskrit \emph{mukha}, \textquotedblleft{}mouth, face\textquotedblright{} (CDIAL 10158), and \ml{കുരു}, \textquotedblleft{}a pustule, a boil\textquotedblright{}, which is native (DEDR 1780).
\end{cousinweb}

\begin{grammarlens}[title={What you've built}]
One more word for describing what you lived and felt.
\end{grammarlens}

\subsection*{Your turn}
\begin{itemize}
\raggedright
  \item \emph{Say it:} \emph{mukhakkuru}
  \item \emph{Say it:} \emph{mukhakkuru}, once more
  \item \emph{Say it:} \emph{aniyattikkŭ mukhattŭ orupāṭŭ mukhakkuru uṇṭŭ.} — My little sister has a lot of pimples on her face.
  \item \emph{Recall:} say the Malayalam for a clinic, then the Malayalam for a walking stick, then say \emph{mukhakkuru} again
\end{itemize}

\begin{tcolorbox}[breakable,colback=teal!4,colframe=teal!35!black,title={Before you move on}]
What does \ml{മുഖക്കുരു} mean? (\textquotedblleft{}A pimple; acne\textquotedblright{}.) Say it once more.
\end{tcolorbox}
\section[\texorpdfstring{\ml{മരവിക്കുക} (maravikkuka)}{മരവിക്കുക (maravikkuka)}]{\ml{മരവിക്കുക} (maravikkuka) — to go numb}

\label{lesson:ML-C540-maravikkuka}



\begin{quote}
Before the new one: say the Malayalam for a walking stick, then the Malayalam for a pimple.
\end{quote}

\subsection*{You'll want to know: \ml{മരവിക്കുക}}
\textbf{\ml{മരവിക്കുക}} — \emph{maravikkuka} — \textquotedblleft{}to go numb\textquotedblright{}.

\begin{quote}
\ml{തണുത്ത} \ml{വെള്ളത്തിൽ} \ml{കൈ} \ml{മരവിച്ചു}.
\end{quote}

\begin{quote}
\emph{taṇutta veḷḷattil kai maraviccu.}
\end{quote}

\begin{quote}
My hand went numb in the cold water.
\end{quote}

Past: \ml{മരവിച്ചു}, \emph{maraviccu}.

Also of feelings: \ml{മനസ്സ്} \ml{മരവിച്ചു}, the heart went numb.

\begin{cousinweb}
\textbf{\ml{മരവിക്കുക}} (\emph{maravikkuka}) looks native, but DEDR has no entry that clearly holds it, so its origin is not settled.
\end{cousinweb}

\begin{grammarlens}[title={What you've built}]
One more word for describing what you lived and felt.
\end{grammarlens}

\subsection*{Your turn}
\begin{itemize}
\raggedright
  \item \emph{Say it:} \emph{maravikkuka}
  \item \emph{Say it:} \emph{maravikkuka}, once more
  \item \emph{Say it:} \emph{taṇutta veḷḷattil kai maraviccu.} — My hand went numb in the cold water.
  \item \emph{Recall:} say the Malayalam for a walking stick, then the Malayalam for a pimple, then say \emph{maravikkuka} again
\end{itemize}

\begin{tcolorbox}[breakable,colback=teal!4,colframe=teal!35!black,title={Before you move on}]
What does \ml{മരവിക്കുക} mean? (\textquotedblleft{}To go numb\textquotedblright{}.) Say it once more.
\end{tcolorbox}
\section[\texorpdfstring{\ml{കഴയ്ക്കുക} (kaḻaykkuka)}{കഴയ്ക്കുക (kaḻaykkuka)}]{\ml{കഴയ്ക്കുക} (kaḻaykkuka) — to feel sore (of limbs)}

\label{lesson:ML-C540-kalaykkuka}



\begin{quote}
Before the new one: say the Malayalam for a pimple, then the Malayalam for to go numb.
\end{quote}

\subsection*{You'll want to know: \ml{കഴയ്ക്കുക}}
\textbf{\ml{കഴയ്ക്കുക}} — \emph{kaḻaykkuka} — \textquotedblleft{}to feel sore (of limbs)\textquotedblright{}.

\begin{quote}
\ml{ഒരുപാട്} \ml{നടന്നതുകൊണ്ട്} \ml{കാലുകൾ} \ml{കഴച്ചു}.
\end{quote}

\begin{quote}
\emph{orupāṭŭ naṭannatukoṇṭŭ kālukaḷ kaḻaccu.}
\end{quote}

\begin{quote}
My legs ached from walking so much.
\end{quote}

Past: \ml{കഴച്ചു}, \emph{kaḻaccu}.

The tired ache in arms and legs after hard use.

\begin{cousinweb}
\textbf{\ml{കഴയ്ക്കുക}} (\emph{kaḻaykkuka}) looks native, but DEDR has no entry that clearly holds it, so its origin is not settled.
\end{cousinweb}

\begin{grammarlens}[title={What you've built}]
One more word for describing what you lived and felt.
\end{grammarlens}

\subsection*{Your turn}
\begin{itemize}
\raggedright
  \item \emph{Say it:} \emph{kaḻaykkuka}
  \item \emph{Say it:} \emph{kaḻaykkuka}, once more
  \item \emph{Say it:} \emph{orupāṭŭ naṭannatukoṇṭŭ kālukaḷ kaḻaccu.} — My legs ached from walking so much.
  \item \emph{Recall:} say the Malayalam for a pimple, then the Malayalam for to go numb, then say \emph{kaḻaykkuka} again
\end{itemize}

\begin{tcolorbox}[breakable,colback=teal!4,colframe=teal!35!black,title={Before you move on}]
What does \ml{കഴയ്ക്കുക} mean? (\textquotedblleft{}To feel sore (of limbs)\textquotedblright{}.) Say it once more.
\end{tcolorbox}
\section[(dialogue)]{Words from chapters 533–540 — a review}

\label{lesson:ML-R540-review-533-540}



\begin{quote}
Say the Malayalam for a vaccine and immunity.
\end{quote}

\subsection*{Your turn}
Say these aloud:

\begin{itemize}
\raggedright
  \item a vaccine, immunity, a prescription, to grow weak, breathlessness
  \item asthma, heartburn, constipation, insomnia, to have a cough
  \item cancer, a limp, a migraine, rheumatism, to throw up
  \item a blood test, a vein, a muscle, the lungs, to take a rest
  \item the brain, the intestines, tears, to snore, to perspire
  \item to get well, to tend, a wheelchair, to hurt, to pass out
  \item to doze off, to pant, to smart, a diet, nutritious food
  \item a clinic, a walking stick, a pimple, to go numb, to feel sore
\end{itemize}

\begin{tcolorbox}[breakable,colback=teal!4,colframe=teal!35!black,title={Before you move on}]
A vaccine? (\textbf{\ml{വാക്സിൻ}}, \emph{vāksin}.) Immunity? (\textbf{\ml{പ്രതിരോധശേഷി}}, \emph{pratirōdhaśēṣi}.) A prescription? (\textbf{\ml{കുറിപ്പടി}}, \emph{kuṟippaṭi}.) To grow weak? (\textbf{\ml{തളരുക}}, \emph{taḷaruka}.) Breathlessness? (\textbf{\ml{ശ്വാസംമുട്ട്}}, \emph{śvāsaṁmuṭṭŭ}.) Asthma? (\textbf{\ml{ആസ്ത്മ}}, \emph{āstma}.) Heartburn? (\textbf{\ml{നെഞ്ചെരിച്ചിൽ}}, \emph{neñcericcil}.) Constipation? (\textbf{\ml{മലബന്ധം}}, \emph{malabandhaṁ}.) Insomnia? (\textbf{\ml{ഉറക്കമില്ലായ്മ}}, \emph{uṟakkamillāyma}.) To have a cough? (\textbf{\ml{ചുമയ്ക്കുക}}, \emph{cumaykkuka}.) Cancer? (\textbf{\ml{അർബുദം}}, \emph{arbudaṁ}.) A limp? (\textbf{\ml{മുടന്ത്}}, \emph{muṭantŭ}.) A migraine? (\textbf{\ml{ചെന്നിക്കുത്ത്}}, \emph{cennikkuttŭ}.) Rheumatism? (\textbf{\ml{വാതം}}, \emph{vātaṁ}.) To throw up? (\textbf{\ml{ഛർദ്ദിക്കുക}}, \emph{charddikkuka}.) A blood test? (\textbf{\ml{രക്തപരിശോധന}}, \emph{raktapariśōdhana}.) A vein? (\textbf{\ml{ഞരമ്പ്}}, \emph{ñarampŭ}.) A muscle? (\textbf{\ml{പേശി}}, \emph{pēśi}.) The lungs? (\textbf{\ml{ശ്വാസകോശം}}, \emph{śvāsakōśaṁ}.) To take a rest? (\textbf{\ml{വിശ്രമിക്കുക}}, \emph{viśramikkuka}.) The brain? (\textbf{\ml{തലച്ചോറ്}}, \emph{talaccōṟŭ}.) The intestines? (\textbf{\ml{കുടൽ}}, \emph{kuṭal}.) Tears? (\textbf{\ml{കണ്ണുനീർ}}, \emph{kaṇṇunīr}.) To snore? (\textbf{\ml{കൂർക്കംവലിക്കുക}}, \emph{kūrkkaṁvalikkuka}.) To perspire? (\textbf{\ml{വിയർക്കുക}}, \emph{viyarkkuka}.) To get well? (\textbf{\ml{സുഖം} \ml{പ്രാപിക്കുക}}, \emph{sukhaṁ prāpikkuka}.) To tend? (\textbf{\ml{പരിചരിക്കുക}}, \emph{paricarikkuka}.) A wheelchair? (\textbf{\ml{വീൽചെയർ}}, \emph{vīlceyar}.) To hurt? (\textbf{\ml{വേദനിക്കുക}}, \emph{vēdanikkuka}.) To pass out? (\textbf{\ml{ബോധം} \ml{കെടുക}}, \emph{bōdhaṁ keṭuka}.) To doze off? (\textbf{\ml{മയങ്ങുക}}, \emph{mayaṅṅuka}.) To pant? (\textbf{\ml{കിതയ്ക്കുക}}, \emph{kitaykkuka}.) To smart? (\textbf{\ml{നീറുക}}, \emph{nīṟuka}.) A diet? (\textbf{\ml{ഭക്ഷണക്രമം}}, \emph{bhakṣaṇakramaṁ}.) Nutritious food? (\textbf{\ml{പോഷകാഹാരം}}, \emph{pōṣakāhāraṁ}.) A clinic? (\textbf{\ml{ക്ലിനിക്ക്}}, \emph{klinikkŭ}.) A walking stick? (\textbf{\ml{ഊന്നുവടി}}, \emph{ūnnuvaṭi}.) A pimple? (\textbf{\ml{മുഖക്കുരു}}, \emph{mukhakkuru}.) To go numb? (\textbf{\ml{മരവിക്കുക}}, \emph{maravikkuka}.) To feel sore? (\textbf{\ml{കഴയ്ക്കുക}}, \emph{kaḻaykkuka}.)
\end{tcolorbox}
